% Generated draft: requires mathematical and source review.
\documentclass[11pt]{article}
\usepackage[margin=25mm]{geometry}
\usepackage{fontspec}
\setmainfont{Latin Modern Roman}
\usepackage{amsmath,amssymb,mathtools,unicode-math}
\setmathfont{Latin Modern Math}
\usepackage{enumitem,needspace,xurl,hyperref}
\setlength{\parindent}{0pt}
\setlength{\parskip}{5pt}
\setlength{\emergencystretch}{3em}
\newcommand{\N}{\mathbb{N}}
\newcommand{\Z}{\mathbb{Z}}
\newcommand{\Q}{\mathbb{Q}}
\newcommand{\R}{\mathbb{R}}
\newcommand{\C}{\mathbb{C}}
\newcommand{\F}{\mathbb{F}}
\newcommand{\A}{\mathbb{A}}
\newcommand{\PP}{\mathbb{P}}
\newcommand{\E}{\mathbb{E}}
\newcommand{\eps}{\varepsilon}
\newcommand{\dd}{\,\mathrm{d}}
\begin{document}
\section*{UnsolvedMath ID: 30004497}
\textbf{Source problem code:} catalogue entry
\subsection{Definitions and mathematical statement}
Let \(\mathbb A_{\mathbb Q}\) denote the ring of adeles of \(\mathbb Q\). The supplied statement uses a symbol \(D\), its group of units \(D^\times\), and the quotient
\[
X=D^\times(\mathbb Q)\backslash D^\times(\mathbb A_{\mathbb Q}).
\]
The record does not define \(D\) or specify hypotheses on it. If \(f\) is a function on \(X\), interpret \(f(y)\), for \(y\in D^\times(\mathbb A_{\mathbb Q})\), as its value at the coset represented by \(y\). The symbols \(\mu\), the Laplace operator, and the Hecke operators are not specified. Formally, a Hecke--Laplace eigenfunction is a simultaneous eigenfunction of the intended Hecke operators and Laplace operator, with \(\lambda\) its Laplace eigenvalue; the operators and their common function space remain to be supplied.

The source asks whether there exists \(\delta>0\) such that, for \(y\in D^\times(\mathbb A_{\mathbb Q})\) and every relevant Hecke--Laplace eigenfunction \(f\) of eigenvalue \(\lambda\),
\[
f(y)=\int_X f(x)\,\mathrm d\mu(x)
       +O\!\left(\lambda^{1/4-\delta}\right).
\]
The single existential quantifier for \(\delta\) precedes the quantifiers over \(y\) and \(f\). The record does not specify whether the implied constant is uniform in \(y\), which eigenfunction family is intended, or which other parameters the constant may depend on. The apparent spectral asymptotic \(\lambda\to+\infty\) is not explicitly stated in the source.

This formula is retained as a proposed assertion with unresolved hypotheses. No normalization of \(f\) is supplied. Since scalar multiples of an eigenfunction remain eigenfunctions, an additive error bound uniform over \(f\) cannot be interpreted without a normalization or another explicit restriction on scaling. The record also does not specify the measure \(\mu\), its normalization, a treatment of the center, or assumptions ensuring convergence of the integral. In particular, the integral cannot be identified with a probability average from the supplied information. These choices must be resolved before this becomes a fully specified mathematical conjecture.
\subsection{Short English statement}
The supplied question asks for a fixed positive power saving in the difference between a Hecke--Laplace eigenfunction's value at an adelic point and its integral over the displayed quotient. The algebra \(D\), eigenfunction normalization, operators, measure, convergence conditions, and error-term uniformity are unspecified, so the intended conjecture requires clarification. The supplied open label is a dated curation claim, not an independently verified mathematical conclusion.
\subsection{Sources}
\begin{itemize}
\item Local UnsolvedMath registry: supplied record with permanent ID \(30004497\), exact source problem code \texttt{catalogue entry}, titled ``Subconvex Pointwise Equidistribution of Hecke--Laplace Eigenfunctions,'' in the Oberwolfach Reports Open Problems collection. No individual proposer is supplied; the proposed year is recorded as 2020.
\item Supplied report citation: \emph{Automorphic Forms and Arithmetic} (2021), Oberwolfach Reports, pp.\ 1317--1383, DOI \texttt{10.4171/owr/2020/26}. Supplied source URL and literature-source link, labeled ``Oberwolfach Report 26/2020'': \url{https://doi.org/10.4171/owr/2020/26}. No specific page or passage containing this question is identified. The background describes these reports as documenting mathematical workshops and their open-problem discussions.
\item Supplied literature assessment, checked on \textbf{2026-08-22}: the record labels the problem open, describes the proposed power-saving estimate for point values and the period integral of all relevant eigenfunctions as a question in the report, and reports that no unconditional resolution of this exact assertion was found. This is supplied dated curation derived from the OWR material, not independently verified here; its time-sensitive status and cited claims require independent verification.
\item The registry labels the statement \texttt{corrected\_verified} and states that its mathematical content was preserved while notation, typography, omitted definitions, or duplicate clauses were repaired for self-contained publication. This supplied editorial designation does not resolve the missing hypotheses in the supplied record or constitute mathematical review of this draft.
\end{itemize}
\end{document}
